\documentclass[11pt,a4paper]{article}
\usepackage[margin=21mm]{geometry}
\usepackage[T1]{fontenc}
\usepackage[utf8]{inputenc}
\usepackage[ngerman]{babel}
\usepackage{booktabs,tabularx,amsmath,siunitx,xcolor,fancyhdr,enumitem}
\definecolor{kitblue}{HTML}{005B82}
\pagestyle{fancy}\fancyhf{}\fancyhead[L]{\small Physikalisches Praktikum}\fancyhead[R]{\thepage}\setlength{\headheight}{14pt}
\begin{document}
\begin{center}
 {\color{kitblue}\LARGE\bfseries Versuchsprotokoll}\par\medskip
 {\large P3/P4 — Dämpfung eines Federpendels}\par
\end{center}
\begin{tabularx}{\linewidth}{@{}>{\bfseries}lX>{\bfseries}lX@{}}
 Name & Vorname Nachname & Datum & 29.08.2026\\
 Versuch & P3/P4 & Gruppe & 07\\
 Betreuer & Dr. Beispiel & Raum & 2.14
\end{tabularx}
\section*{Aufgabenstellung}
Bestimmen Sie die Eigenfrequenz und die Dämpfung eines Federpendels aus einer
Zeitreihe der Auslenkung. Vergleichen Sie die Messwerte mit dem Modell eines
gedämpften harmonischen Oszillators.
\section*{Aufbau und Durchführung}
Feder, Masse, Wegsensor und Datenaufnahme werden auf einer stabilen Schiene
montiert. Das Pendel wird um eine definierte Strecke ausgelenkt und ohne
zusätzlichen Anstoß freigegeben. Die Auslenkung wird mit \SI{100}{Hz}
aufgezeichnet; der Versuch wird dreimal wiederholt.
\section*{Messwerte}
\begin{tabular}{@{}l@{\quad}rr@{\quad}r@{}}
\toprule
Messung & $T$ (s) & $x_0$ (cm) & $x_5$ (cm)\\\midrule
1 & 1.24 & 4.80 & 2.92\\
2 & 1.25 & 4.75 & 2.88\\
3 & 1.23 & 4.82 & 2.95\\
\bottomrule
\end{tabular}
\section*{Auswertung}
Für $x(t)=x_0 e^{-\delta t}\cos(\omega t+\varphi)$ ergibt sich aus den
aufeinanderfolgenden Maxima beispielhaft
\[T=(1.24\pm0.02)\ \si{\second},\qquad
\delta=(0.10\pm0.01)\ \si{\per\second}.\]
\section*{Fehlerbetrachtung und Fazit}
Diskutieren Sie Sensorauflösung, Ablesezeitpunkt, Reibung und die Annahme
einer linearen Feder. Die drei Periodendauern stimmen innerhalb der
Messunsicherheit überein; das Modell beschreibt die Messreihe angemessen.
\section*{Unterschriften}
\noindent Student: \rule{45mm}{.4pt}\hfill Betreuer: \rule{40mm}{.4pt}
\end{document}
